\documentclass[12pt,a4paper]{article}
\usepackage[utf8]{inputenc}
\usepackage[english]{babel}
\usepackage{amsmath, amssymb, amsthm}
\usepackage{geometry}
\geometry{margin=2.5cm}

\usepackage{kantlipsum}

\usepackage{booktabs}

\usepackage[shortlabels]{enumitem}

\usepackage{mathtools}

\theoremstyle{plain}

\newtheorem*{ejercicio1}{Ejercicio 2.8}
\newtheorem*{ejercicio2}{Ejercicio 2.11}
\newtheorem*{ejercicio3}{Ejercicio 2.16}
\newtheorem*{ejercicio4}{Ejercicio 2.17}
\newtheorem*{ejercicio5}{Ejercicio 2.40}
\newtheorem*{ejercicio6}{Ejercicio 3.2}
\newtheorem*{ejercicio7}{Ejercicio 3.17}
\newtheorem*{ejercicio27}{Ejercicio 2.7}
\newtheorem*{af}{Afirmación}


\newcommand{\A}{\mathbb{A}}
\newcommand{\C}{\mathbb{C}}
\newcommand{\R}{\mathbb{R}}
\newcommand{\kk}{k}
\newcommand{\Ga}{\Gamma}
\newcommand{\OO}{\mathcal{O}}
\newcommand{\Ker}{\mathrm{Ker}\,}

\title{Tarea 4. Introducción a la Geometría Algebraica MAT2824}
\author{Felipe López}

\begin{document}
\maketitle

En el ejercicio 2.8 utilizaremos el ejercicio 2.7, por lo que este último será enunciado primero.

\begin{ejercicio27}
	Si $\varphi \colon V \to W$ es un mapa polinomial y $X$ es un subconjunto algebraico irreducible de $W$, demuestre que $\varphi^{-1}(X)$ es un subconjunto algebraico de $V$. Si $\varphi^{-1}(X)$ es irreducible y $X$ está contenido en la imagen de $\varphi$, demuestre que $X$ es irreducible.
\end{ejercicio27}

\begin{ejercicio1}
	\begin{enumerate}[(a)]
		\item Demuestre que $\{(t,t^{2},t^{3}) \in \mathbb{A}^{3}(k) \ | \ t \in k\}$ es una variedad afin.

		\item Demuestre que $V(xz - y^{2}, yz - x^{3}, z^{2} - x^{2}y) \subset \mathbb{A}^{3}(\C)$ es una variedad. (Hint: $y^{3} - x^{4}, z^{3} - x^{5}, z^{4} - y^{5} \in I(V)$. Encuentre un mapa polinomial de $\mathbb{A}^{1}(\C)$ a $V$ sobreyectivo.)
	\end{enumerate}
\end{ejercicio1}

\begin{proof}
	\boxed{a.} Por un lado, sabemos que $\{(t,t^2,t^3)\in\A^3(k): t\in k\}=:V$ es un conjunto algebraico por T1. Luego, nos falta demostrar irreducibilidad en $\A^3(k)$. Para ello, recordamos que $V=V(y-x^2,z-x^3)$ e $I(V(I))= I$ para todo ideal $I \subset k[x,y,z]$. Por lo tanto, como $V(\mathcal{F})=V((\mathcal{F}))$,
	\[
	  I(V(y-x^2,z-x^3))=(y-x^2, z-x^3).
	\]
	Luego, sea $\Psi \colon k[x,y,z] \to k[t]$ tal que $f(x,y,z)=f(t,t^2,t^3)$. Así, es claro que $\Psi$ es sobreyectivo y que $\ker\Psi=I(V)$, pues el kernel consiste de todos los polinomios $f\in k[x,y,z]$ tales que $f(t,t^2,t^3)=0$, y eso es exactamente $I(V)$. Luego, por el primer Teorema de Isomorfismos:
	\[
	  \frac{k[x,y,z]}{I(V)}\;\simeq\; k[t].
	\]
	$\therefore$ $k[x,y,z]/I(V)$ dominio entero $\Rightarrow$ $I(V)$ ideal primo $\Rightarrow$ $V(I(V))=V$ irreducible (que se cumple pues $k$ es alg. cerrado).

	\bigskip
	\boxed{b.} Sean $V:=V(xz-y^2,\,yz-x^3,\,z^2-x^2y)\subset\A^3(\C)$. Por el hint, tenemos que $y^3-x^4,\, z^3-x^5,\, z^4-y^5\in I(V)$.

	Luego, notamos que, para todo $(x,y,z)\in V$ se cumple que $y^3=x^4$, $z^3=x^5$, $z^4=y^5$. Es decir,
	\begin{align*}
		\left(\frac{y}{x}\right)^3 = x \; &\Rightarrow\; y = x\cdot\frac{y}{x} = y = \left(\frac{y}{x}\right)^4 \\
		&\Rightarrow\; z^3 = \left(\left(\frac{y}{x}\right)^3\right)^5 = \left(\frac{y}{x}\right)^{15} \\
		&\Rightarrow\; z = \left(\frac{y}{x}\right)^5.
	\end{align*}
	Por lo tanto, si $t:=\dfrac{y}{x}$, entonces $x=t^3$, $y=t^4$, $z=t^5$. Esto nos da un claro candidato para la construcción del mapa polinomial. Luego, definimos $\Psi:\A^1(\C)\to V$ como $t\mapsto(t^3,t^4,t^5)$. \par

	Notamos que si $f_1=xz-y^2$, $f_2=yz-x^3$, $f_3=z^2-x^2y$, entonces
	\begin{align*}
	  f_1(t^3,t^4,t^5) &= t^3\cdot t^5-(t^4)^2 = t^8-t^8=0,\\
	  f_2(t^3,t^4,t^5) &= t^4 t^5-(t^3)^3 = t^9-t^9=0,\\
	  f_3(t^3,t^4,t^5) &= (t^5)^2-(t^3)^2 t^4 = t^{10}-t^{10}=0.
	\end{align*}
	Por lo tanto, $\Psi$ está bien definida.

	Luego, notamos que es sobreyectiva, pues para $(x,y,z)\in V$ basta con tomar $t=\tfrac{y}{x}$. En efecto, $x=t^3$, $y=t^4$, $z=t^5$ se cumple por lo visto gracias al hint.

	Por último, por el ejercicio 2.7 del Fulton, tenemos que como $\Psi^{-1}(X) = \A^{1}(\C)$ es irreducible y $V \subseteq \Psi(\A^{1}(\C) = V$, entonces $V$ es irreducible.
\end{proof}

%-----------------------------------------------------------------------

\begin{ejercicio2}	
	Sean $f \in \Gamma(V)$, $V$ una variedad. Definamos
	\[ G(f) = \{(a_{1}, \dots, a_{n}, a_{n + 1}) \in \mathbb{A}^{1} \ | \ (a_{1},\dots,a_{n}) \in V \text{ y } a_{n + 1} = f(a_{1},\dots,a_{n})\}, \]
	el grafico de $f$. Demuestre que $G(f)$ es una variedad afin y que el mapa $(a_{1},\dots,a_{n}) \mapsto (a_{1},\dots,a_{n},f(a_{1},\dots,a_{n}))$ define un isomorfismo de $V$ con $G(f)$. (la proyección da la inversa.)
\end{ejercicio2}

\begin{proof}
	Queremos ver que $G(f)$ es un conjunto algebraico y que es también irreducible.

	Para ver lo primero, sea
	\[
	  J=(I(V),\,x_{n+1}-f(x_1,\ldots,x_n)) \subset k[x_1,\ldots,x_{n+1}].
	\]
	Así, notar que:
	\begin{align*}
		V(J) &= \{(a_1,\ldots,a_n,a_{n+1}) \in \A^{n+1} : g(a_1,\ldots,a_n) = 0\;\forall g \in I(V) \;\wedge\; a_{n+1}-f(a_1,\ldots,a_n)=0\bigr\} \\
		&= G(f).
	\end{align*}
	Es decir, $G(f)$ es un conjunto algebraico.

	Para ver irreducibilidad, veamos que $J$ es ideal primo. En efecto, sea $\Psi:k[x_1,\ldots,x_{n+1}]\to\Gamma(V)=k[x_1,\ldots,x_n]/I(V)$ tal que $x_i\mapsto\bar{x}_i$ $\forall i\leq n$ y $x_{n+1}\mapsto f(x_1,\ldots,x_n)$. El mapa es claramente sobreyectivo, pues $\Gamma(V)$ está generado por los $x_i$ ($i\leq n$) y $f\in\Gamma(V)$. Luego, afirmamos que $\Ker\Psi=J$.

	\boxed{\supseteq} Sea $g\in J$. Entonces, $g=aF+b(x_{n+1}-f)$ para algunos $a,b\in k[x_1,\ldots,x_{n+1}]$ y $F \in I(V)$. Luego, notar que $aF=0$ en $\Gamma(V)$, pues $F \in I(V)$, y como $x_{n+1}\mapsto f \Rightarrow x_{n+1}-f=0 \Rightarrow g \mapsto 0\Rightarrow g\in\Ker\Psi\Rightarrow\Ker\Psi\supseteq J$.

	\boxed{\subseteq} Sea $G\in\Ker\Psi$. Notar que $x_{n+1}-f\in k[x_1,\ldots,x_{n+1}]$ es un polinomio mónico, por lo que podemos dividir por él. Es decir, podemos escribir
	\[
	  G = p(x_{n+1}-f)+r(x_1,\ldots,x_n)
	\]
	para algún $r\in k[x_1,\ldots,x_n]$ y $p\in k[x_1,\ldots,x_{n+1}]$.

	Luego, notar que
	\begin{align*}
		\Psi(G) &= \bar{p}(f-f)+\bar{r}\\
		&= 0+\bar{r} = 0
	\end{align*}
	$\Rightarrow\bar{r}=0$. Por lo tanto, $\bar{r}\in I(V)$
	$\Rightarrow G=p(x_{n+1}-f)+r\in J$
	$\Rightarrow\Ker\Psi\subseteq J$. $\therefore\Ker\Psi=J$.

	Luego, por el primer Teorema de Isomorfismos:
	\[
	  \frac{k[x_1,\ldots,x_{n+1}]}{J}\;\simeq\;\Gamma(V)
	\]
	$\Rightarrow$ $J$ ideal primo, pues $\Gamma(V)$ es dominio entero (pues $V$ es irreducible)
	$\Rightarrow$ $V(J)=G(f)$ irreducible
	$\Rightarrow$ $G(f)$ variedad afín.

	Por otro lado, sea $\alpha:V\to G(f)$ tal que $(a_1,\ldots,a_n)\mapsto(a_1,\ldots,a_n,f(a_1,\ldots,a_n))$ y sea $\pi:G(f)\to V$ tal que $(a_1,\ldots,a_n,a_{n+1})\mapsto(a_1,\ldots,a_n)$. Notar que como la identidad y $f$ son polinomios, entonces ambos, $\alpha$ y $\pi$, son mapas polinomiales.

	Luego, notar que si $(a_1,\ldots,a_n,a_{n+1})\in G(f)$, entonces $a_{n+1}=f(a_1,\ldots,a_n)$ y
	\begin{align*}
	  \alpha\circ\pi(a_1,\ldots,a_{n+1}) &= \alpha(a_1,\ldots,a_n)\\
	  &=(a_1,\ldots,a_n,f(a_1,\ldots,a_n))\\
	  &=(a_1,\ldots,a_n,a_{n+1}).
	\end{align*}
	Similarmente, si $(a_1,\ldots,a_n)\in V$, entonces:
	\begin{align*}
	  \pi\circ\alpha(a_1,\ldots,a_n) &= \pi(a_1,\ldots,a_n,f(a_1,\ldots,a_n))\\
	  &=(a_1,\ldots,a_n).
	\end{align*}
	Es decir, $\alpha$ es un isomorfismo entre $V$ y $G(f)$.
\end{proof}

%-----------------------------------------------------------------------
\begin{ejercicio3}
	Sea $k = \C$. Demos a $\mathbb{A}^{n}(\C) = \C^{n}$ la topología usual (obtenida por identificar a $\C$ con $\R^{2}$ y por lo tanto, a $\C^{n}$ con $\R^{2n}$). Recordar que un espacio topológico es \textit{conexo por caminos} si para cualquier $P,Q \in X$, existe un mapa continuo $\gamma \colon [0,1] \to X$ tal que $\gamma(0) = P, \gamma(1) = Q$.
	\begin{enumerate}[(a)]
		\item Demuestre que $\C \setminus S$ es conexo por caminos para cualquier conjunto finito $S$.

		\item Sea $V$ un conjunto algebraico en $\mathbb{A}^{n}(\C)$. Demuestre que $\mathbb{A}^{n}(\C) \setminus V$ es conexo por caminos. (Hint: si $P,Q \in \mathbb{A}^{n}(\C) \setminus V$, sea $L$ la linea que pasa por $P$ y $Q$. Entonces, $L \cap V$ es finito y $L$ es isomorfo a $\mathbb{A}^{1}(\C)$.)
	\end{enumerate}
\end{ejercicio3}

\begin{proof}
	\boxed{a.} Dado que $\C$ es Hausdorff, si $|S|=d$, podemos encontrar $d$ bolas abiertas disjuntas tales que cada bola contenga a un punto de $S$. Luego, sean $P,Q\in\C\setminus S$. Tomamos la línea recta que conecta a $P$ y $Q$. Si esta línea pasa por uno de los puntos de $S$, podemos desviar la recta por el borde de la Bola que contiene a tal punto, formando un semicírculo. Esto es claramente continuo y conecta a los puntos $P$ y $Q$. (Podemos asegurarnos de tomar las bolas tales que $P$ y $Q$ no esté contenida en ninguna).

	\bigskip
	\boxed{b.} Sean $P,Q\in\A^n(\C)\setminus V$ tales que $P=(a_1,\ldots,a_n)$ y $Q=(b_1,\ldots,b_n)$, y sea
	\[
	  L:=\bigl\{(a_1+t(b_1-a_1),\ldots,a_n+t(b_n-a_n)):t\in\C\bigr\}=\{P+t(Q-P):t\in\C\}
	\]
	la recta que pasa por $P$ y $Q$. Por el hint, sabemos que $L\cap V$ es finita y que $L$ es isomorfa a $\A^1(\C)=\C$ como variedades.

	Notamos que $V=V(I)$ para algún ideal $I\subset\C[x_1,\ldots,x_n]$. Luego, sea $\Psi(t)=P+t(Q-P)$ la parametrización de la recta $L$. Si $f\in I$, entonces $f\circ\Psi$ sigue siendo un polinomio (pues $\Psi$ es lineal), pero en variable $t$. Luego,
	\[
	  \Psi(t)\in V \iff f(\Psi(t))=0\quad\forall f\in I.
	\]
	Definimos
	\[
	  T:=\{t\in\C\mid (f\circ\Psi)(t)=0\;\forall f\in I\}.
	\]
	Entonces,
	\[
	  T = V\bigl(\{f\circ\Psi\in\C[t]:f\in I\}\bigr)\subseteq\C
	\]
	$\Rightarrow$ $T$ conjunto algebraico en $\C$. Como $\C$ es un campo $\Rightarrow\C[t]$ es un DIP $\Rightarrow\bigl(\{f\circ\Psi\in\C[t]\mid f\in I\}\bigr)=(G)$ para algún $G$ $\Rightarrow T=V(G)$.

	Notar que $P,Q\notin V$, por lo que $t=0$ y $t=1$ no pertenecen a $T$ $\Rightarrow G\not\equiv 0\Rightarrow G$ tiene finitas raíces $\Rightarrow T$ es finito. Luego, $L\cap V=\Psi(T)$ también es finito, pues $\Psi$ es inyectiva.

	Como $L\cap V=\Psi(T)$ y $\Psi$ define un homeomorfismo de $L$ con $\C$ $\Rightarrow L\setminus(L\cap V)\simeq\C\setminus T$. Por parte (a), $\C\setminus T$ es conexo por caminos, por lo que $L\setminus(L\cap V)$ es también conexo por caminos.

	Por lo tanto, $\exists\,\delta:[0,1]\to\C\setminus T$ tal que $\delta(0)=0$ y $\delta(1)=1$. Luego, $\delta\circ\Psi:[0,1]\to L\setminus(L\cap V)\subset\A^n(\C)\setminus V$ tal que $\delta\circ\Psi(0)=P$, $\delta\circ\Psi(1)=Q$.

	$\Rightarrow\A^n(\C)\setminus V$ es conexo por caminos.
\end{proof}

%-----------------------------------------------------------------------

\begin{ejercicio4}
	Sea $V = V(y^{2} - x^{2}(x + 1)) \subset \mathbb{A}^{2}$ y sean $\overline{x}, \overline{y}$ los residuos de $x,y$ en $\Gamma(V)$. Además, sea $z = \overline{y} \slash \overline{x} \in k(V)$. Encuentre el conjunto de polos de $z$ y de $z^{2}$.
\end{ejercicio4}

\begin{proof}
	Para ver el caso de $z^2$, notamos que
	\begin{align*}
		& \bar{y}^2 = \bar{y}^2 - (\bar{y}^2-\bar{x}^2(\bar{x}+1)) = \bar{x}^2(\bar{x}+1) \\
		\Rightarrow \ & z^2 = \frac{\bar{y}^2}{\bar{x}^2} = \frac{\bar{x}^2(\bar{x}+1)}{\bar{x}^2} = \bar{x}+1.
	\end{align*}
	Luego, como $k[x,y]\subseteq\OO_p(V)$ para todo $p$, entonces $z^2$ está definido para todo $p\in V$. $\Rightarrow$ El conjunto de polos es vacío.

	Para ver el de $z$, notar que en esta representación de $z$, los puntos donde $\bar{x}$ se anula son donde $x=0$. Luego, notar que $y^2-0^2\cdot(0+1)=0\iff y^2=0\iff y=0$. Por lo tanto, el conjunto de polos de $z$ sólo puede estar formado por el punto $(0,0)$. Verifiquemos que este punto es, en efecto, un polo.

	Para ello, supongamos que no lo es. Es decir, $z\in\OO_{(0,0)}(V)$, por lo que $\exists\,a,b\in k[x,y]$ tal que $z=\dfrac{a}{b}$ con $b(0,0)\neq 0$. Luego, notar que:
	\begin{align*}
		\frac{\bar{y}}{\bar{x}} = \frac{\bar{a}}{\bar{b}} &\Rightarrow\; \bar{y}\bar{b}=\bar{x}\bar{a}\in\Gamma(V) \\
		&\Rightarrow\; yb-xa=0\in(y^2-x^2(x+1)) \\
		&\Rightarrow\; \exists\,c\in k[x,y]:\; yb-xa=c(x,y)(y^2-x^2(x+1)).
	\end{align*}
	Dado que $b(0,0)\neq 0$, el polinomio $b$ tiene un término constante $r\neq 0$. Es decir, el polinomio $yb-xa$ cuenta con un término de grado 1 $ry$. Pero notamos que el polinomio $c(x,y)(y^2-x^2-x^3)$ no puede tener términos de grado menor a 2. Pero esto sería una contradicción, puesto que implicaría que $b(0,0)=0$.

	En conclusión, $z\notin\OO_{(0,0)}(V)$ $\Rightarrow$ El conjunto de polos de $z$ es $\{(0,0)\}$.
\end{proof}

%-----------------------------------------------------------------------

\begin{ejercicio5}
	\begin{enumerate}[(a)]
		\item Supongamos que $I,J$ son ideales comaximales en $R$. Demuestre que $I + J^{2} = R$. Demuestre también que $I^{m}$ y $J^{m}$ son comaximales para todo $m,n$.

		\item Supongamos que $I_{1},\dots,I_{N}$ son ideales de $R$ y que $I_{i}$ y $J_{i} = \bigcap_{j \neq i} I_{J}$ son comaximales para todo $i$. Demuestre que $I_{1}^{n} \cap \cdots \cap I_{N}^{n} = (I_{1} \cdots I_{N})^{n} = (I_{1} \cap \cdots \cap I_{N})^{n}$ para todo $n$.
	\end{enumerate}
\end{ejercicio5}

\begin{proof}
	\boxed{a.1.} Notemos que $\exists\,a\in I$, $b\in J$ tales que $a+b=1\Rightarrow(a+b)^2=1^2\Rightarrow a^2+2ab+b^2=1$. Luego, como $a(a+2b)\in I$ y $b^2\in J^2$, entonces $I+J^2=R$.

	\medskip
	\boxed{a.2.} Notar que existen $a\in I$, $b\in J$ tales que $a+b=1$. Entonces,
	\begin{align*}
	(a+b)^{m+n} &= \sum_{k=0}^{m+n}\binom{m+n}{k}a^{m+n-k}b^k\\
	&= \sum_{k=0}^{n}\binom{m+n}{k}a^{m+n-k}b^k + \sum_{k=n+1}^{m+n}\binom{m+n}{k}a^{m+n-k}b^k.
	\end{align*}
	Notar que en $\sum_{k=0}^{m}\binom{m+n}{k}a^{m+n-k}b^k$ se cumple que $m+n-k\geq m$ $\forall k\in\{0,1,\ldots,n\}$. Es decir:
	\[
	  \sum_{k=0}^{m}\binom{m+n}{k}a^{m+n-k}b^k = a^m\left(\sum_{k=0}^{m}\binom{m+n}{k}a^{n-k}b^k\right)
	\]
	y además, $\left(\sum_{k=0}^{m}\binom{m+n}{k}a^{n-k}b^k\right)\in R$ y $a^m\in I^m$.

	Similarmente, en $\sum_{k=m+1}^{m+n}\binom{m+n}{k}a^{m+n-k}b^k$, se cumple que $k\geq n$ $\forall k\in\{n+1,\ldots,m+n\}$, por lo que:
	\[
	  \sum_{k=m+1}^{m+n}\binom{m+n}{k}a^{m+n-k}b^k = b^n\left(\sum_{k=m+1}^{m+n}\binom{m+n}{k}a^{m+n-k}b^{k-n}\right).
	\]
	Además, $\left(\sum_{k=m+1}^{m+n}\binom{m+n}{k}a^{m+n-k}b^{k-n}\right)\in R$ y $b^n\in J^n$.

	Luego,
	\[
	  \left(\sum_{k=0}^{m}\binom{m+n}{k}a^{m+n-k}b^k\right)\in I^m\quad\text{y}\quad\left(\sum_{k=m+1}^{m+n}\binom{m+n}{k}a^{m+n-k}b^k\right)\in J^n
	\]
	y como
	\[
	  \sum_{k=0}^{m}\binom{m+n}{k}a^{m+n-k}b^k + \sum_{k=m+1}^{m+n}\binom{m+n}{k}a^{m+n-k}b^k = (a+b)^{m+n}=1^{m+n}=1
	\]
	Entonces, $I^m+J^n=R \;\Rightarrow\; I^m+J^n \text{ son coprimos}$.
\end{proof}

%-----------------------------------------------------------------------

\begin{ejercicio6}
	Encuentre los puntos multiple y las lineas tangentes en esos puntos multiple para cada una de las siguientes curvas:
	\begin{itemize}[(a)]
	  \item $y^3-y^2+x^3-x^2+3xy^2+3x^2y+2xy$
	  \item $x^4+y^4-x^2y^2$
	  \item $x^3+y^3-3x^2-3y^2+3xy+1$
	  \item $y^2+(x^2-5)(4x^4-20x^2+25)$
	\end{itemize}
\end{ejercicio6}

\begin{proof}
	\boxed{a.} Sea $F:=y^3-y^2+x^3-x^2+3xy^2+3x^2y+2xy$. Entonces:
	\[
	  \begin{cases}
		F_y(x,y)=3y^2-2y+6xy+3x^2+2x,\\
		F_x(x,y)=3x^2-2x+3y^2+6xy+2y.
	  \end{cases}
	\]
	\[
	  \Rightarrow\begin{cases}
		F_x(x,y)=3(x+y)^2-2(y-x),\\
		F_y(x,y)=3(x+y)^2+2(x-y).
	  \end{cases}
	\]
	Notar que los puntos múltiples $(a,b)$ son aquellos que satisfacen $F_x(a,b)=F_y(a,b)=F(a,b)=0$. Luego, tomamos el sistema de ecuaciones:
	\[
	  \begin{cases}
		(1)\; F_x(x,y)=0\\
		(2)\; F_y(x,y)=0
	  \end{cases}
	  \;\xrightarrow{(1)-(2)}\; 4(x-y)=0\;\Rightarrow\; x=y.
	\]
	$\Rightarrow F_x(x,y)=F_x(x,x)=3(2x)^2=12x^2=0\;\Rightarrow\; x=0\;\Rightarrow\; y=0$. Luego, notamos que $F(0,0)=0$, por lo que $(0,0)$ es el único punto múltiple.

	Dado que el punto es el $(0,0)$, no es necesario trasladar la curva. Luego, notar que
	\begin{align*}
	  F(x,y) &= y^3-y^2+x^3-x^2+3xy^2+3x^2y+2xy\\
			  &= (x+y)^3+(y-x)^2.
	\end{align*}
	Es decir, la recta $x=y$ tiene multiplicidad 2, siendo así la única recta tangente del punto múltiple $(0,0)$.

	\bigskip
	\boxed{b.} Sea $F=x^4+y^4-x^2y^2$. Notemos que:
	\[
	  F_x(x,y)=4x^3-2xy^2,\qquad F_y(x,y)=4y^3-2x^2y.
	\]
	\[
	  \Rightarrow\begin{cases}
		(1)\; 4x^3-2xy^2=0\;\Rightarrow\; 2x(2x^2-y^2)=0,\\
		(2)\; 4y^3-2x^2y=0\;\Rightarrow\; 2y(2y^2-x^2)=0.
	  \end{cases}
	\]
	Por $F_x=0$: $x=0$ ó $2x^2=y^2$. \par
	En el primer caso, $x=0\Rightarrow 2y(2y^2)=0\Rightarrow y=0$. \par
	En el segundo caso, $2x^2=y^2\Rightarrow 2y(4x^2-x^2)=0\Rightarrow 2y(3x^2)=0\Rightarrow 3y(2x^2)=0\Rightarrow 3y^3=0\Rightarrow y=0$.

	Como $F(0,0)=0$ $\Rightarrow(0,0)$ es el único punto múltiple.

	Luego, podemos tratar a $F=y^4-x^2y^2+x^4$ como una ecuación cuadrática en $y^2$. Es decir,
	\begin{align*}
	  y^2 &= \frac{x^2\pm\sqrt{x^4-4\cdot 1\cdot x^4}}{2} = \frac{x^2\pm\sqrt{x^4-4x^4}}{2}\\
		  &= \frac{x^2\pm x^2\sqrt{-3}}{2} = x^2\!\left(\frac{1\pm i\sqrt{3}}{2}\right).
	\end{align*}
	Para simplificar notación, notar que
	\[
	  e^{i\pi/3}=\cos\!\left(\tfrac{\pi}{3}\right)+i\sin\!\left(\tfrac{\pi}{3}\right)=\frac{1}{2}+i\frac{\sqrt{3}}{2},
	\]
	\[
	  e^{-i\pi/3}=\cos\!\left(\tfrac{\pi}{3}\right)-i\sin\!\left(\tfrac{\pi}{3}\right)=\frac{1}{2}-i\frac{\sqrt{3}}{2}.
	\]
	$\Rightarrow y^2=x^2 e^{\pm i\pi/3}\;\Rightarrow\; y=\pm x\,e^{\pm i\pi/6}$. \par
	Así, hemos obtenido 4 rectas tangentes distintas $y=\pm xe^{i\pi/6}$ e $y=\pm xe^{-i\pi/6}$.

	\bigskip
	\boxed{c.} Sea $F:=x^3+y^3-3x^2-3y^2+3xy+1$. Notar que:
	\[
	  F_x(x,y)=3x^2-6x+3y=0 \text{ y } F_y(x,y)=3y^2-6y+3x=0.
	\]
	Luego,
	\begin{align*}
		y=2x-x^2\;&\Rightarrow\;(2x-x^2)^2-2(2x-x^2)+x=0 \\
		&\Rightarrow 4x^2-4x^3+x^4-4x+2x^2+x=0 \\
		&\Rightarrow\; x^4-4x^3+6x^2-3x=0 \\
		&\Rightarrow x(x^3-4x^2+6x-3)=0 \\
		&\Rightarrow\; x(x-1)(x^2-3x+3)=0
	.\end{align*}
	Notar que las soluciones de $x^2-3x+3=0$ están dadas por
	\[
	  \frac{3\pm\sqrt{9-12}}{2}=\frac{3\pm i\sqrt{3}}{2}=\frac{3}{2}\pm i\frac{\sqrt{3}}{2}.
	\]
	Entonces
	\[
	  F_x(x,y)=x(x-1)\!\left(x-\tfrac{3}{2}-i\tfrac{\sqrt{3}}{2}\right)\!\left(x-\tfrac{3}{2}+i\tfrac{\sqrt{3}}{2}\right).
	\]
	Por lo tanto, $x=0,\;1,\;\tfrac{3}{2}+i\tfrac{\sqrt{3}}{2},\;\tfrac{3}{2}-i\tfrac{\sqrt{3}}{2}$. Vamos por casos: \par

	Si $x=0\Rightarrow 3y=0\Rightarrow y=0$. Luego, $F(0,0)=1\neq 0$. \par

	Si $x=1\Rightarrow -3+3y=0\Rightarrow y=1$. Luego, $F(1,1)=1+1-3-3+3+1=0$.\par

	Si $x=\dfrac{3}{2}+i\dfrac{\sqrt{3}}{2}$, entonces $y=-\left(\dfrac{3}{2}+i\dfrac{\sqrt{3}}{2}\right)^2+2\left(\dfrac{3}{2}+i\dfrac{\sqrt{3}}{2}\right)$:

	\begin{align*}
	  y &= -\left(\frac{9}{4}+2i\frac{3\sqrt{3}}{4}-\frac{3}{4}\right)+3+i\sqrt{3}\\
		&= -\frac{3}{2}-i\frac{3\sqrt{3}}{2}+3+i\sqrt{3}\\
		&= \frac{3}{2}-i\frac{\sqrt{3}}{2}.
	\end{align*}

	\par Similiarmente, si $x=\dfrac{3}{2}-i\dfrac{\sqrt{3}}{2}\Rightarrow y=\dfrac{3}{2}+i\dfrac{\sqrt{3}}{2}$.

	Luego, notar que $F(0,0)=1\neq 0$. Además,
	\begin{align*}
	  F\!\left(\tfrac{3}{2}+i\tfrac{\sqrt{3}}{2},\,\tfrac{3}{2}-i\tfrac{\sqrt{3}}{2}\right)
	  &= (-3)\!\left(\tfrac{3}{2}+i\tfrac{\sqrt{3}}{2}\right)+(-3)\!\left(\tfrac{3}{2}-i\tfrac{\sqrt{3}}{2}\right)
		+3\!\left(\tfrac{3}{2}+i\tfrac{\sqrt{3}}{2}\right)\!\left(\tfrac{3}{2}-i\tfrac{\sqrt{3}}{2}\right)+1\\
	  &= -\tfrac{9}{2}-i\tfrac{3\sqrt{3}}{2}-\tfrac{9}{2}+i\tfrac{3\sqrt{3}}{2}
		+3\!\left(\tfrac{9}{4}+\tfrac{3}{4}\right)+1\\
	  &= -\tfrac{18}{2}+\tfrac{36}{4}+1 = 1\neq 0.
	\end{align*}
	Similarmente, $F\!\left(\tfrac{3}{2}-i\tfrac{\sqrt{3}}{2},\,\tfrac{3}{2}+i\tfrac{\sqrt{3}}{2}\right)=1\neq 0$.

	Luego, el único punto múltiple es $(1,1)$. Debemos hacer un cambio de coordenadas. Para ello, sean $u=x-1$ y $v=y-1$ (i.e., $x=u+1$, $y=v+1$):
	\begin{align*}
	  F(x,y) &= F(u+1,v+1)\\
	  &= (u+1)^3+(v+1)^3-3(u+1)^2-3(v+1)^2+3(u+1)(v+1)+1\\
	  &= (u^3+3u^2+3u+1)+(v^3+3v^2+3v+1)-3(u^2+2u+1)\\
	  &\quad -3(v^2+2v+1)+3(uv+u+v+1)+1\\
	  &= u^3+v^3+3uv.
	\end{align*}
	El término de menor grado es $3uv$. Es decir, tenemos dos rectas tangentes distintas: $u=0$, $v=0$. Si volvemos a las coordenadas originales obtenemos las rectas $x=1$ e $y=1$.

	\bigskip
	\boxed{d.} Sea $F:=y^2+(x^2-5)(4x^4-20x^2+25)=y^2+(x^2-5)(2x^2-5)^2$. Notar que $F_y = 2y$, por lo que $y = 0$, y
	\begin{align*}
	  F_x &= (2x)(2x^2-5)^2+(x^2-5)\cdot 2(2x^2-5)(4x)\\
		  &= (2x)(2x^2-5)^2+8x(x^2-5)(2x^2-5)\\
		  &= 2x(2x^2-5)\bigl(2x^2-5+4(x^2-5)\bigr)\\
		  &= 2x(2x^2-5)(2x^2-5+4x^2-20)\\
		  &= 2x(2x^2-5)(6x^2-25),
	\end{align*}

	Observar que un punto múltiple debe cumplir $F(x,0)=0\Rightarrow(x^2-5)(2x^2-5)^2=0\Rightarrow x^2=5$ ó $x^2=\tfrac{5}{2}$. Si comparamos con $F_x\Rightarrow x=0$ ó $x^2=\dfrac{5}{2}$ ó $x^2=\dfrac{25}{6}$.

	Es decir, los puntos múltiples son
	\[
	  \left(\sqrt{\tfrac{5}{2}},\,0\right)\quad\text{y}\quad\left(-\sqrt{\tfrac{5}{2}},\,0\right).
	\]
	Luego, hacemos un cambio de coordenadas $u=x-\sqrt{5/2}$, $v=y$:
	\begin{align*}
	  F(u+\sqrt{5/2},v) &= v^2+\bigl((u+\sqrt{5/2})^2-5\bigr)\bigl(2(u+\sqrt{5/2})^2-5\bigr)^2\\
	  &= v^2+\bigl(u^2+\sqrt{10}\,u+\tfrac{5}{2}-5\bigr)\bigl(2(u^2+\sqrt{10}\,u+\tfrac{5}{2})-5\bigr)^2\\
	  &= v^2+\bigl(u^2+\sqrt{10}\,u-\tfrac{5}{2}\bigr)\bigl(2u^2+2\sqrt{10}\,u\bigr)^2\\
	  &= v^2+u^2\bigl(u^2+\sqrt{10}\,u-\tfrac{5}{2}\bigr)\bigl(2u+2\sqrt{10}\bigr)^2\\
	  &= v^2+u^2\bigl(u^2+\sqrt{10}\,u-\tfrac{5}{2}\bigr)(4u^2+4\sqrt{10}\,u+40).
	\end{align*}
	Notamos que los términos de menor grado serán aquellos que resulten en grado 2. Luego,
	\[
	  F_2 = v^2-\tfrac{5}{2}\cdot 40u^2 = v^2-100u^2=(v-10u)(v+10u).
	\]
	Así, hemos obtenido las dos rectas $v = 10u$ y $v = -10u$ tangentes. Si volvemos a las coordenadas originales, obtenemos $y = 10(x - \sqrt{5 \slash 2})$ e $y = -10(x - \sqrt{5 \slash 2})$, que son las rectas tangentes al punto múltiple $(\sqrt{5 \slash 2}, 0)$. \par
	El procedimiento es similar para encontrar las rectas tangentes al punto múltiples $(-\sqrt{5 \slash 2}, 0)$.
\end{proof}

\end{document}
